\chapter{Nokta Tahmini; Parametre, Tahmin Edici ve Tahmin}
\label{ch:point-estimation}

\begin{prismbox}[PrismLight]{Öğrenme çıktıları}
Bu bölümün sonunda:
\begin{itemize}
  \item parametre, tahmin edici ve gerçekleşmiş tahmini ayırt edebilecek,
  \item ortalama, oran, varyans ve farklar için nokta tahmini hesaplayabilecek,
  \item tahmin edicileri yanlılık, varyans, MSE ve tutarlılık bakımından karşılaştırabilecek,
  \item bir nokta tahminini uygun hassasiyet ve belirsizlik bilgisiyle raporlayabileceksiniz.
\end{itemize}
\end{prismbox}

\begin{prerequisitebox}
Parametre ve istatistik ayrımı, örnekleme dağılımı, beklenen değer, varyans
ve standart hata kavramları bilinmelidir.
\end{prerequisitebox}

Parametre--istatistik ayrımı Bölüm~\ref{sec:population-parameter}'de,
örnekleme dağılımı ve standart hata ise Bölüm~\ref{ch:sampling-distributions}
ile Bölüm~\ref{sec:standard-error}'da geliştirilmiştir.

\section{Parametre, tahmin edici ve tahmin}
\label{sec:estimator-estimate}

\begin{definitionbox}
Parametre evrene ait sabit fakat bilinmeyen bir niceliktir. Tahmin edici,
örneklemden hesaplanan rassal bir kuraldır; tahmin ise gözlenen örneklem için
bu kuralın aldığı sayısal değerdir. Örneğin $\bar X$ evren ortalaması $\mu$
için bir tahmin edici, $\bar x=72{,}4$ ise elde edilen tahmindir.
\end{definitionbox}

\section{Sık kullanılan nokta tahminleri}

\begin{table}[H]
\centering
\caption{Parametreler ve doğal örneklem tahminleri.}
\label{tab:natural-estimators}
\begin{tabular}{p{0.38\textwidth}p{0.36\textwidth}}
\toprule
Hedef parametre & Nokta tahmini \\
\midrule
Evren ortalaması $\mu$ & Örneklem ortalaması $\bar x$ \\
Evren oranı $p$ & Örneklem oranı $\hat p=x/n$ \\
Evren varyansı $\sigma^2$ & Örneklem varyansı $s^2$ \\
İki evren ortalaması farkı & $\bar x_1-\bar x_2$ \\
Korelasyon $\rho$ & Örneklem korelasyonu $r$ \\
\bottomrule
\end{tabular}
\end{table}

Tek bir nokta tahmini belirsizliği göstermez. $\bar x=72$ sonucu, örneklem
hacmi 10 olduğunda ve 1000 olduğunda aynı görünür; oysa tahminlerin hassasiyeti
çok farklıdır. Bu nedenle nokta tahmini standart hata veya güven aralığıyla
birlikte verilmelidir.

\section{Yanlılık}
\label{sec:estimator-bias}

Bir tahmin edici için
\[
\text{Yanlılık}(\hat\theta)
=\E[\hat\theta]-\theta
\]
tanımlanır. Beklenen değeri hedef parametreye eşitse tahmin edici yansızdır.
Yansızlık uzun dönemli bir özelliktir; tek bir örneklemde tahminin parametreye
eşit olacağını söylemez.

\section{Varyans, standart hata ve ortalama karesel hata}
\label{sec:estimator-mse}

Tahmin edicinin örneklemden örnekleme değişkenliği varyansıyla ölçülür.
Standart hata bu varyansın kareköküdür. Yanlılık ile değişkenliği birleştiren
ortalama karesel hata
\[
\MSE(\hat\theta)=\Var(\hat\theta)
+\text{Yanlılık}(\hat\theta)^2
\]
eşitliğini sağlar.

\begin{sezgibox}
Hedef tahtasında atışların merkezi gerçek parametreyi temsil etsin. Atışların
merkezden sistematik kayması yanlılık, birbirinden uzaklığı varyanstır. İyi bir
tahmin edici hem hedefe yakın hem de tekrarlarda kararlı olmalıdır.
\end{sezgibox}

Şekil~\ref{fig:estimator-bias-spread}, aynı hedef parametre için dört
şematik örnekleme dağılımını gösterir. Soldaki dağılımların merkezi hedefte,
sağdakilerin merkezi hedefin sağındadır. Üst sıradaki dar eğriler daha küçük
varyansı gösterir. Dolayısıyla dar bir dağılım tek başına yansızlık anlamına
gelmez; geniş bir yansız dağılım da tek örneklemde hedefe yakınlığı garanti etmez.

\begin{figure}[H]
\centering
\begin{tikzpicture}[font=\footnotesize]
  \foreach \sx/\sy/\center/\spread/\heading in {
    -2.7/0/0/0.4/{Yansız, düşük varyans},
    2.7/0/0.85/0.4/{Yanlı, düşük varyans},
    -2.7/-3.2/0/0.85/{Yansız, yüksek varyans},
    2.7/-3.2/0.85/0.85/{Yanlı, yüksek varyans}}{
    \begin{scope}[shift={(\sx,\sy)}]
      \node[font=\small\bfseries] at (0,2.1) {\heading};
      \draw[->,PrismGray] (-2.2,0)--(2.3,0);
      \draw[PrismBlue,dashed,line width=0.7pt] (0,0)--(0,1.65);
      \draw[PrismBlue,line width=1pt,domain=-2.15:2.15,samples=90,smooth]
        plot (\x,{0.56/\spread*exp(-((\x-\center)^2)/(2*\spread^2))});
      \node[below] at (0,0) {$\theta$};
      \node[below,align=center] at (0,-0.4) {Tahmin değeri};
    \end{scope}
  }
\end{tikzpicture}
\caption{Yanlılık ve varyansın ayrı özellikler olarak okunması. Kesikli
dikey çizgi gerçek parametreyi, eğriler tekrarlı örneklemlerdeki tahmin
dağılımlarını gösterir; çizimler şematiktir.}
\label{fig:estimator-bias-spread}
\end{figure}

\section{Tutarlılık}

Örneklem hacmi büyüdükçe tahmin edici gerçek parametre çevresinde yoğunlaşıyorsa
tutarlı olarak adlandırılır. Örneklem ortalaması, sonlu varyans altında evren
ortalaması için tutarlıdır. Tutarlılık büyük örneklem özelliğidir ve küçük
örneklem performansını tek başına garanti etmez.

\section{Etkinlik ve sağlamlık}

Aynı parametreyi tahmin eden iki yansız tahmin ediciden varyansı daha küçük
olan, daha \textbf{etkin} kabul edilir. Ancak düşük varyans her koşulda en iyi
uygulamalı tercih anlamına gelmez. Örneğin örneklem ortalaması normal veya
yaklaşık simetrik verilerde etkinken, medyan ağır kuyruklu ve aykırı değerli
dağılımlarda daha \textbf{sağlam} olabilir.

Bu nedenle tahmin edici seçimi yalnız tek bir matematiksel özelliğe değil,
hedef parametreye, veri üretim sürecine ve olası model sapmalarına dayanır.
Yansızlık, düşük varyans, küçük MSE ve sağlamlık bazen aynı tahmin edicide
birleşirken bazen aralarında bir denge kurulması gerekir.
Bu ölçütlerin karar kuramı ve tahmin kuramı içindeki ilişkisi için
\citet{hogg2015} temel bir başvuru kaynağıdır.

\section{Raporlama hassasiyeti}

Nokta tahmini, ölçme aracının hassasiyetini ve tahmin belirsizliğini aşan
ondalık basamaklarla raporlanmamalıdır. Standart hatası 2,4 puan olan bir
ortalamanın 72,4387 biçiminde verilmesi gerçek bilgi düzeyini olduğundan
yüksek gösterir. Genellikle tahmin ile standart hata aynı ondalık basamakta
sunulur ve kullanılan ağırlıklandırma ya da örnekleme tasarımı belirtilir.

\section{Bir tahmini değerlendirirken}

\begin{enumerate}
  \item Hedef parametre açık mı?
  \item Örnekleme tasarımı hedef evrene uygun mu?
  \item Tahminin standart hatası veya güven aralığı verilmiş mi?
  \item Aykırı değerler ya da eksik veriler tahmini belirgin etkiliyor mu?
  \item Tahmin bağlama uygun birim ve hassasiyetle raporlanmış mı?
\end{enumerate}

\section{Çözümlü örnek: aynı tahmin, farklı hassasiyet}

İki araştırmada olumlu yanıt oranı $\hat p=0{,}60$ bulunsun. İlk araştırmada
$n=25$, ikincisinde $n=400$ kişi vardır. Büyük örneklem yaklaşımıyla
\[
\SE(\hat p)=\sqrt{\frac{\hat p(1-\hat p)}{n}}
\]
olduğundan standart hatalar sırasıyla yaklaşık $0{,}098$ ve $0{,}024$'tür.
Nokta tahminleri aynı olsa da ikinci sonuç çok daha hassastır. Bununla
birlikte ikinci örneklem kolayda örnekleme ile seçilmişse düşük standart hata
temsil yanlılığı sorununu ortadan kaldırmaz.

\begin{outputguidebox}
Bir tahmin tablosunda önce hedef parametreyi ve birimi belirleyin. Ardından
örneklem hacmini, nokta tahminini, standart hatayı ve varsa güven aralığını
birlikte okuyun. Gereğinden fazla ondalık basamak sahte hassasiyet yaratır.
\end{outputguidebox}

\section{Yazılım uygulaması: tahmin ve standart hata}

\begin{softwarebox}
\begin{lstlisting}
import numpy as np

responses = np.array([1, 0, 1, 1, 0, 1, 1, 0, 1, 0])
p_hat = responses.mean()
se = np.sqrt(p_hat * (1 - p_hat) / len(responses))

print("Oran tahmini:", p_hat)
print("Yaklasik standart hata:", se)

rng = np.random.default_rng(2026)
simulated = rng.binomial(1, .40, size=(10000, 50))
estimates = simulated.mean(axis=1)
print("Tahminlerin merkezi:", estimates.mean())
print("Ampirik standart hata:", estimates.std(ddof=1))
\end{lstlisting}
İkili kodlamada 1 olayın gerçekleşmesini, 0 gerçekleşmemesini gösterdiğinde
aritmetik ortalama örneklem oranına eşittir. Benzetimde tahminlerin merkezi
0,40'a, standart sapması ise kuramsal $\sqrt{0{,}40(0{,}60)/50}$ değerine
yaklaşmalıdır.
\end{softwarebox}

\begin{mistakebox}
Bir tahminin yansız olması her örneklemde doğru değeri vereceği anlamına
gelmez. Yansızlık, aynı süreç çok kez tekrarlandığında tahminlerin merkezine
ilişkin bir özelliktir.
\end{mistakebox}

\begin{assumptionbox}
Standart hata hesabı kullanılan tahmin ediciye ve örnekleme tasarımına
bağlıdır. Ağırlıklı, kümeli veya bağımlı verilerde basit rastgele örneklem
formülleri belirsizliği olduğundan küçük gösterebilir.
\end{assumptionbox}

\section{Çözümlü problemler}

\subsection*{Problem 1: Parametre, tahmin edici ve tahmin}

Bir üniversitedeki öğrencilerin çevrim içi derslere düzenli katılım oranı
araştırılmaktadır. Rastgele seçilen 80 öğrencinin 54'ü düzenli katılmıştır.
Parametreyi, tahmin ediciyi ve gerçekleşmiş tahmini belirleyin.

\textbf{Çözüm.} Üniversitedeki bütün öğrencilerin düzenli katılım oranı
bilinmeyen parametre $p$'dir. $X$ örneklemde düzenli katılan öğrenci sayısı
olmak üzere
\[
\hat p=\frac{X}{n}
\]
örneklem oranı bir tahmin edicidir; örneklem seçilmeden önce rassaldır.
Gözlenen örnekteki gerçekleşmiş tahmin
\[
\hat p=\frac{54}{80}=0{,}675
\]
olur. Buna göre örneklemde düzenli katılım oranı yüzde 67,5'tir. Bu değer
evren oranının kesin değeri değil, örneklemden elde edilen nokta tahminidir.

\textbf{Raporlama örneği.} ``Rastgele seçilen 80 öğrencinin 54'ü düzenli
katılım göstermiş; evren oranı için nokta tahmini yüzde 67,5 bulunmuştur.''

\subsection*{Problem 2: İki yansız tahmin edicinin karşılaştırılması}

$X_1$ ve $X_2$, ortalaması $\mu$ ve varyansı $\sigma^2$ olan bir evrenden
bağımsız iki gözlem olsun. $T_1=X_1$ ve
$T_2=(X_1+X_2)/2$ tahmin edicilerini yansızlık ve varyans bakımından
karşılaştırın.

\textbf{Çözüm.} Beklenen değerlerin doğrusallığından
\[
\E[T_1]=\E[X_1]=\mu
\]
ve
\[
\E[T_2]=\frac{\E[X_1]+\E[X_2]}{2}
=\frac{\mu+\mu}{2}=\mu
\]
olur. Her iki tahmin edici de yansızdır. Varyansları ise
\[
\Var(T_1)=\sigma^2,
\qquad
\Var(T_2)=\frac{\sigma^2+\sigma^2}{4}
=\frac{\sigma^2}{2}
\]
biçimindedir. $T_2$ aynı hedef için yarı varyansa sahip olduğundan daha
etkindir. İkinci gözlemi kullanmak tahminin örneklemden örnekleme
değişkenliğini azaltmıştır.

\textbf{Yorum.} İki tahmin edicinin de yansız olması onları eşit kalitede
yapmaz; değişkenlik de karşılaştırılmalıdır.

\subsection*{Problem 3: MSE ile yanlılık--varyans dengesi}

Aynı $\theta$ parametresi için A tahmin edicisinin yanlılığı 0 ve varyansı 9;
B tahmin edicisinin yanlılığı 2 ve varyansı 3 olsun. Her iki tahmin edicinin
MSE değerini hesaplayıp kare hata ölçütüne göre karşılaştırın.

\textbf{Çözüm.} A tahmin edicisi için
\[
\MSE(A)=9+0^2=9.
\]
B tahmin edicisi için
\[
\MSE(B)=3+2^2=7.
\]
A yansız olmasına rağmen yüksek varyansı nedeniyle MSE değeri daha büyüktür.
B iki birim yanlıdır; ancak değişkenliğindeki azalma yanlılığın karesinden
daha ağır bastığı için toplam karesel hata 7'ye düşmüştür. Kare hata kaybı
altında B tercih edilir.

Bu sonuç ``yanlı tahmin edici her zaman daha iyidir'' anlamına gelmez. Tercih,
yanlılık ve varyansın büyüklüklerine, örneklem hacmine ve yanlış tahminlerin
uygulamadaki maliyetine bağlıdır.

\subsection*{Problem 4: Örneklem varyansında neden $n-1$ kullanılır?}

$2,4,6$ gözlemleri için örneklem ortalamasını, kareli sapmalar toplamını ve
$n-1$ bölenli örneklem varyansını hesaplayın. Aynı toplam $n$'ye bölünürse
hangi değer elde edilir?

\textbf{Çözüm.} Örneklem ortalaması
\[
\bar x=\frac{2+4+6}{3}=4
\]
ve kareli sapmalar toplamı
\[
(2-4)^2+(4-4)^2+(6-4)^2=4+0+4=8
\]
olur. Örneklem varyansı
\[
s^2=\frac{8}{3-1}=4
\]
olarak hesaplanır. Aynı toplam $n=3$'e bölünürse $8/3\approx2{,}67$ elde
edilir.

Sapmalar örneklem ortalamasına göre hesaplandığında bir serbestlik derecesi
kullanılmış olur. $n-1$ düzeltmesi, bağımsız ve aynı dağılımlı örneklemede
$s^2$'yi evren varyansı $\sigma^2$ için yansız tahmin edici yapar. Tek bir
örnekte 4 değerinin gerçek evren varyansına eşit olması gerekmez; yansızlık
tekrarlı örnekleme özelliğidir.

\subsection*{Problem 5: Alt gruplardan birleşik ortalama tahmini}

Bir araştırmada 40 lisans öğrencisinin ortalama çalışma süresi 70 dakika,
60 lisansüstü öğrencinin ortalaması 80 dakika bulunmuştur. Bütün 100
öğrencinin birleşik ortalamasını hesaplayın. Grup ortalamalarının basit
ortalamasını almak neden yanlış olur?

\textbf{Çözüm.} Birleşik toplam çalışma süresi
\[
40(70)+60(80)=2800+4800=7600
\]
dakikadır. Bu nedenle birleşik ortalama
\[
\bar x=\frac{7600}{100}=76
\]
dakika olur. Grup ortalamalarının basit ortalaması $(70+80)/2=75$ dakika
olurdu; ancak bu hesap iki grubun eşit büyüklükte olduğunu varsayar. Burada
lisansüstü grup daha büyük olduğu için 80 dakikalık ortalama birleşik tahminde
daha fazla ağırlık taşır.

\textbf{Varsayım kontrolü.} Bu birleşik ortalama gözlenen 100 kişi için
doğrudur. Gruplar evrende farklı oranlarda bulunuyor veya farklı seçilme
olasılıklarıyla örneklenmişse evren ortalaması için tasarım ağırlıkları ayrıca
kullanılmalıdır.

\begin{outputguidebox}
Bir tahmin çıktısında tahmin edilen büyüklüğün ortalama, oran, fark veya başka
bir parametre olup olmadığını belirleyin. Nokta tahminiyle birlikte standart
hata, örneklem hacmi, kullanılan ağırlıklar ve yuvarlama düzeyini kontrol
edin. Küçük standart hata, kötü örnekleme tasarımından gelen yanlılığı ortadan
kaldırmaz.
\end{outputguidebox}

\section{Literatür verisiyle yazılım uygulamaları}
\label{sec:spss-b07}

\textbf{Amaç ve veri.} \texttt{b07.csv} üzerinden ortalama ve oran için
nokta tahmini üretin \citep{cortez2008data}. \texttt{pass10}, yıl sonu notu
10 veya üzerinde ise 1'dir.

\textbf{Ortak veri.} Doğrulanan B07 uygulamasında üç yazılım da
\nolinkurl{companion/spss/csv/b07.csv} dosyasını kullanmıştır. Dosya
395 satır ve 10 sütun içerir. B06'nın seçim göstergeleri ve ağırlığı bu
uygulamada kullanılmaz; 395 kaydın tamamı analiz edilir. Bu karşılaştırma
Excel içe aktarmanın doğrulandığı anlamına gelmez. Kodları
\texttt{companion} klasörünü içeren paket kökünden çalıştırın;
veri sözlüğü ve hazırlık için bkz. sayfa~\pageref{sec:spss-guide}.

\subsection{IBM SPSS uygulaması}
\textbf{Başlamadan önce.} Aşağıdaki komut dosyası B07 CSV verisini
açar ve özetler. Menü yolunu kullanacaksanız önce veri açma ve tanımlama
komutlarını \texttt{EXECUTE} satırı dahil çalıştırın. Filtre, ağırlık ve
dosya bölme kapalı olmalıdır. Komut yolundaki \texttt{AGGREGATE}, etkin
veriyi tek özet satırıyla değiştirir; bu özeti ham veri dosyasının üzerine
kaydetmeyin. Analizi tekrarlamak için kodu veri açma adımından başlatın.

\begin{spssadimlar}
\spssadim{\emph{Analyze $\to$ Descriptive Statistics $\to$ Frequencies} yolunda \texttt{pass10} seçin ve tabloyu alın. 1 kategorisinin oranını belirleyin.}
\spssadim{\emph{Data $\to$ Aggregate} yolunu açın. \emph{Break Variable(s)} alanını boş bırakın; böylece dosyanın tamamı tek grup olarak özetlenir.}
\spssadim{\texttt{g3} için \emph{Mean} ve \emph{Standard deviation}, \texttt{pass10} için \emph{Mean} işlevlerini seçin. Özetlere sırasıyla \texttt{ort}, \texttt{ss}, \texttt{oran} adlarını verin; grup gözlem sayısını \texttt{n} adıyla ekleyin.}
\spssadim{Yalnız özet değişkenleri içeren yeni bir veri kümesi oluşturmayı seçin. \emph{OK} sonrasında o veri kümesini etkinleştirin; bir satır bütün dosyanın özetidir.}
\spssadim{\emph{Transform $\to$ Compute Variable} yolunda sırayla \texttt{sh\_ort=ss/SQRT(n)} ve \texttt{sh\_oran=SQRT(oran*(1-oran)/n)} hesaplarını yapın. Sol tarafı hedef değişken, sağ tarafı sayısal ifade alanına girin.}
\spssadim{\emph{Data View} içinde özetleri inceleyin. Bir çıktı listesi için aşağıdaki tam komut dosyasını baştan çalıştırın; bu yol ham veriyi yeniden açar ve özet veriyi etkin dosyanın yerine koyar.}
\end{spssadimlar}

{\raggedright
\noindent\textbf{Komutların görevi.} \texttt{AGGREGATE}, boş \texttt{BREAK} ile bütün dosyayı tek özet satırına indirir. \texttt{COMPUTE} standart hataları hesaplar; \texttt{FORMATS} görünümü, \texttt{LIST} çıktı listesini düzenler.\par}

\textbf{Uygulama kodu:} \texttt{b07}. Doğrulanan veri dosyası
\texttt{companion/spss} altındaki \texttt{csv/b07.csv} dosyasıdır.

\begin{lstlisting}[style=prismSPSS]
SET DECIMAL=DOT.
GET DATA /TYPE=TXT
 /FILE='companion/spss/csv/b07.csv'
 /ENCODING='UTF8'
 /ARRANGEMENT=DELIMITED /DELCASE=LINE
 /DELIMITERS="," /QUALIFIER='"'
 /FIRSTCASE=2
 /VARIABLES=id F8.0 school F8.0 sex F8.0 age F8.0
 studytime F8.0 absences F8.0 g1 F8.0 g2 F8.0
 g3 F8.0 pass10 F8.0.
FILTER OFF.
WEIGHT OFF.
SPLIT FILE OFF.
VARIABLE LABELS
 g3 'Yil sonu matematik notu'
 pass10 'Yil sonu notu 10 veya uzerinde'.
VALUE LABELS pass10
 0 '10 puanin altinda' 1 '10 puan ve uzerinde'.
VARIABLE LEVEL
 id school sex pass10 (NOMINAL)
 studytime (ORDINAL)
 age absences g1 g2 g3 (SCALE).
EXECUTE.
FREQUENCIES VARIABLES=pass10.
AGGREGATE /OUTFILE=*
 /BREAK= /n=N /ort=MEAN(g3) /ss=SD(g3)
 /oran=MEAN(pass10).
COMPUTE sh_ort=ss/SQRT(n).
COMPUTE sh_oran=SQRT(oran*(1-oran)/n).
FORMATS n (F8.0).
FORMATS ort ss oran sh_ort sh_oran (F14.9).
LIST VARIABLES=n ort ss sh_ort oran sh_oran.
\end{lstlisting}

\begin{outputguidebox}
\textbf{Paylaşılan SPSS çıktısıyla doğrulanan değerler.}
\texttt{LIST} çıktısında $n=\SPContextN$, ortalama \SPContextMean,
$s=\SPContextSD$, ortalamanın standart hatası \SPContextSE\ bulunmuştur.
Eşiği karşılayan \SPContextPass\ öğrenci vardır; oran \SPContextPHat,
bağımsız Bernoulli modeli altındaki oran standart hatası \SPContextPSE'dır.
Ortalama ve oran farklı tahmin hedefleridir; örneklem oranı evren oranının
kesin değeri değildir. \texttt{AGGREGATE} veriyi tek özet satırına indirir;
çıktıdaki bir satır yalnız bir öğrencinin ölçüldüğü anlamına gelmez.
\end{outputguidebox}


\par\noindent\begin{minipage}{\linewidth}
\subsection{R uygulaması}
\label{sec:r-gercek-b07}

İkili \texttt{pass10} değişkeninin ortalaması oran tahminidir. Ortalama ve oran standart hataları farklı formüllerle hesaplanır.

\begin{lstlisting}[style=prismR]
veri <- read.csv("companion/spss/csv/b07.csv")
stopifnot(nrow(veri) == 395, ncol(veri) == 10,
          !anyNA(veri),
          all(veri$pass10 == as.integer(veri$g3 >= 10)))
gozlem_sayisi <- nrow(veri)
ortalama <- mean(veri$g3)
standart_sapma <- sd(veri$g3)
oran <- mean(veri$pass10)
frekans <- table(veri$pass10)
print(frekans); print(100 * prop.table(frekans))
print(c(n=gozlem_sayisi, ort=ortalama, ss=standart_sapma,
  sh_ort=standart_sapma/sqrt(gozlem_sayisi), oran=oran,
  sh_oran=sqrt(oran*(1-oran)/gozlem_sayisi)))
print(sum(veri$g3 == 0))
\end{lstlisting}

\textbf{Çıktıyı okuma.} Paylaşılan R çıktısı 395 satır, 10 sütun ve
her sütunda sıfır eksik değer gösterir. \texttt{pass10} ile
\texttt{g3 >= 10} koşulu bütün kayıtlarda uyumludur. \texttt{ort} ve
\texttt{oran} nokta tahminleri; \texttt{sh\_ort} ve \texttt{sh\_oran}
ise bunların kullanılan model altındaki standart hatalarıdır.
Sıfır notlu 38 kayıt analizde tutulmuştur.
\end{minipage}\par

\par\noindent\begin{minipage}{\linewidth}
\subsection{Python uygulaması}
\label{sec:python-b07}

\texttt{ddof=1} örneklem standart sapmasını seçer; \texttt{sh\_ort} ile \texttt{sh\_oran} ayrı belirsizlik ölçüleridir.

\begin{lstlisting}[style=prismPythonApp]
import pandas as pd
import numpy as np

veri = pd.read_csv("companion/spss/csv/b07.csv")
assert veri.shape == (395, 10)
assert not veri.isna().any().any()
assert veri.pass10.eq((veri.g3 >= 10).astype(int)).all()
gozlem_sayisi = len(veri)
ortalama = veri.g3.mean()
standart_sapma = veri.g3.std(ddof=1)
oran = veri.pass10.mean()
frekans = veri.pass10.value_counts().sort_index()
print(frekans); print(100 * frekans / gozlem_sayisi)
print(pd.Series({
    "n": gozlem_sayisi, "ort": ortalama, "ss": standart_sapma,
    "sh_ort": standart_sapma / np.sqrt(gozlem_sayisi),
    "oran": oran,
    "sh_oran": np.sqrt(oran*(1-oran)/gozlem_sayisi)
}))
print(int((veri.g3 == 0).sum()))
\end{lstlisting}

\textbf{Çıktıyı okuma.} Paylaşılan Python çıktısında da 395 satır,
10 sütun, sıfır eksik değer ve eşik kodlaması için tam uyum vardır.
Altı özet değer R ile eşleşir; 130 ve 265 kişilik frekanslar ile
38 sıfır notlu kayıt da doğrulanmıştır. Bu uygulama bir hipotez testi,
$p$ değeri veya güven aralığı üretmez. R ve Python ana veriyi korur.
\end{minipage}\par

\subsection{Sonuçların karşılaştırılması ve ortak rapor}
12 Eylül 2026'da paylaşılan SPSS tabloları ile R ve Python metin
çıktıları karşılaştırılmıştır. Gözlem sayısı, ortalama, standart sapma,
ortalamanın standart hatası, oran ve oranın standart hatasından oluşan
altı değer R ve Python'da gösterilen 13 ondalık basamakta eşleşmiştir.
SPSS listesindeki ondalıklı değerler dokuz basamakta aynı sonuçlarla
uyumludur. Frekanslar aynıdır; yüzdeler gösterilen yuvarlama
hassasiyetlerinde uyumludur. Proje CSV'sinden yapılan hesaplar da
bu sonuçları destekler. Aşağıdaki tablolar çıktıların kitap düzenine
uyarlanmış özetidir; ekran görüntüsü değildir.

\begin{table}[H]
\centering
\caption{B07'de yıl sonu notunun 10 puan eşiğine göre dağılımı}
\label{tab:b07-dogrulanmis-frekans}
\begin{tabular}{@{}lrr@{}}
\toprule
Kategori & Frekans & Yüzde \\
\midrule
10 puanın altında & 130 & 32,911392 \\
10 puan ve üzerinde & 265 & 67,088608 \\
Toplam & 395 & 100,000000 \\
\bottomrule
\end{tabular}
\par\smallskip
{\footnotesize Yüzdeler altı ondalık basamağa yuvarlanmıştır.
Eksik değer olmadığından SPSS'in \emph{Percent} ve \emph{Valid Percent}
sütunları eşittir; yüzdelerin paydası 395 kayıttır.\par}
\end{table}

\begin{table}[H]
\centering
\caption{B07'nin doğrulanmış nokta tahminleri ve standart hataları}
\label{tab:b07-dogrulanmis-tahmin}
\begin{tabular}{@{}lr@{}}
\toprule
Ölçü & Değer \\
\midrule
Gözlem sayısı & 395 \\
Ortalama not & 10,415190 \\
Notların standart sapması & 4,581443 \\
Ortalamanın standart hatası & 0,230517 \\
10 puan eşiğini karşılayanların oranı & 0,670886 \\
Oranın standart hatası & 0,023643 \\
\bottomrule
\end{tabular}
\par\smallskip
{\footnotesize Ondalıklı değerler altı basamağa yuvarlanmıştır.
Standart sapma 394 böleniyle hesaplanmıştır. Standart hatalar
$s/\sqrt{395}$ ve $\sqrt{\hat p(1-\hat p)/395}$ formüllerine dayanır;
ara hesaplarda yuvarlanmamış değerler kullanılmıştır.\par}
\end{table}

\Needspace{15\baselineskip}
\begin{outputguidebox}
Tablo~\ref{tab:b07-dogrulanmis-frekans} içindeki 265 kişi,
$\hat p=265/395$ oranını verir. Tablo~\ref{tab:b07-dogrulanmis-tahmin}
içindeki 4,581443, bireysel notların yayılımını; 0,230517 ise kullanılan
bağımsız gözlemler modeli altında ortalama tahmininin standart hatasını
gösterir. Oranın standart hatası 0,023643, yaklaşık 2,3643 yüzde puanıdır;
\%0,023643 olarak okunmaz.

Bu standart hatalar öğretim amacıyla kullanılan modelin hesaplarıdır;
kaynak araştırmanın örnekleme tasarımına uygun belirsizlik hesabının
veya daha geniş bir evrene genellemenin doğrulandığı anlamına gelmez.
Yalnız eldeki 395 kaydın betimlenmesi hedefleniyorsa ortalama ve oran
bu dosyanın özetleridir. Üç yazılımın uyumu temsil gücünü kanıtlamaz.

SPSS'in \texttt{LIST} çıktısındaki bir satır, 395 öğrencinin toplu
özetidir. Ham dosyada 38 sıfır not vardır; bunlar eksik değer sayılmamış,
hem ortalama hesabında hem de eşik altı kategorisinde korunmuştur.
\end{outputguidebox}

\Needspace{9\baselineskip}
\textbf{Raporlama örneği (doğrulanmış çıktılarla).}
``395 kayıtta ortalama yıl sonu notu 10,415190, standart sapması
4,581443 bulunmuştur. Notu 10 veya üzerinde olan 265 kişinin oranı
0,670886'dır (\%67,088608). Bağımsız gözlemler modeli altında
ortalamanın standart hatası 0,230517, oranın standart hatası 0,023643
olarak hesaplanmıştır. Sıfır notlu 38 kayıt analizde tutulmuştur.
Sonuçlar kaynak dosyayı özetlemektedir; daha geniş bir öğrenci evrenine
genellenebilirlik iddiası taşımamaktadır.''


\textbf{Görev.} Nokta tahminlerinin yanına hedef evren tanımını ekleyin.
Sıfır notları gerekçesiz biçimde eksik saymanın analiz paydasını ve
tahmin hedefini nasıl değiştireceğini açıklayın. Yalnız \texttt{g3}
değişkenini değiştirmek ile aynı kayıtları \texttt{pass10} hesabından da
çıkarmak arasındaki farkı belirtin; 99 yazmanın tek başına eksik değer
tanımlamak olmadığını unutmayın. SPSS'te sonraki uygulamaya geçmeden
önce tek satırlık özet yerine gerekli ham veri dosyasını açın.

\section{R ile çözümlü uygulama}
\label{sec:r-b07}

\textbf{Soru.} Aşağıdaki on ikili yanıt için oran tahminini ve yaklaşık
standart hatayı bulun. Ardından gerçek oranı 0,40 olan Bernoulli evreninde
$n=50$ için tahmin edicinin yanlılığını ve MSE'sini benzetimle inceleyin.

\begin{lstlisting}[style=prismR]
yanit <- c(1, 0, 1, 1, 0, 1, 1, 0, 1, 0)
oran <- mean(yanit)
c(oran = oran,
  standart_hata = sqrt(oran * (1 - oran) / length(yanit)))
set.seed(2026)
gercek_oran <- 0.40
hacim <- 50
tahminler <- rbinom(10000, size = hacim,
                    prob = gercek_oran) / hacim
c(merkez = mean(tahminler),
  yanlilik = mean(tahminler) - gercek_oran,
  ampirik_se = sd(tahminler),
  mse = mean((tahminler - gercek_oran)^2))
c(kuramsal_se = sqrt(0.40 * 0.60 / hacim),
  kuramsal_mse = 0.40 * 0.60 / hacim)
hacimler <- c(25, 400)
sqrt(0.60 * 0.40 / hacimler)
\end{lstlisting}

\begin{outputguidebox}
\textbf{Çözüm ve kontrol değerleri.} On yanıtta altı olumlu sonuç vardır:
$\hat p=0{,}60$ ve yaklaşık $\widehat{\SE}=0{,}1549$.
Benzetimin merkezi yaklaşık 0,40, yanlılığı yaklaşık sıfır olmalıdır.
Kuramsal standart hata 0,06928, MSE 0,0048'dir. $\hat p=0{,}60$ için
$n=25$ ve $n=400$ kullanılırsa yaklaşık standart hatalar 0,09798 ve 0,02449 olur.
\end{outputguidebox}

\textbf{Yorum.} İlk hesap bilinmeyen oran yerine örneklem oranını koyar;
benzetimin kuramsal kontrolü ise bilinen 0,40 oranını kullanır. Sonlu sayıdaki
tekrarda görülen küçük ampirik yanlılık, tahmin edicinin kuramsal olarak
yanlı olduğunu kanıtlamaz. On gözlem için normal yaklaşım temkinle kullanılmalıdır.

\textbf{Pekiştirme ve yanıt.} Bernoulli örneklem oranı yansız olduğundan
kuramsal MSE varyansa eşittir: $p(1-p)/n=0{,}0048$.

\section{Bölüm özeti}

\begin{itemize}
  \item Parametre hedef büyüklük, tahmin edici rassal kural, tahmin gözlenen sayısal sonuçtur.
  \item İyi bir tahmin edici yanlılık ve değişkenlik bakımından birlikte değerlendirilir.
  \item MSE, varyans ile yanlılığın karesini aynı ölçüde birleştirir.
  \item Nokta tahmini standart hata veya güven aralığı olmadan hassasiyeti göstermez.
  \item Aynı parametreyi tahmin eden yansız tahmin ediciler varyanslarına göre
  etkinlik bakımından karşılaştırılabilir.
  \item Birleşik tahminlerde alt grup büyüklükleri ve seçilme olasılıkları
  uygun ağırlıklarla hesaba katılmalıdır.
\end{itemize}

\section*{Alıştırmalar}

\begin{enumerate}
  \item Tahmin edici ile tahmin arasındaki farkı açıklayın.
  \item Yansız fakat yüksek varyanslı bir tahmin edicinin sorununu belirtin.
  \item Nokta tahmininin neden tek başına yeterli olmadığını açıklayın.
  \item $\hat p=0{,}40$ ve $n=100$ için yaklaşık standart hatayı hesaplayın.
  \item Düşük standart hatanın neden düşük yanlılık garantisi olmadığını açıklayın.
  \item 120 öğrencinin 78'inin olumlu yanıt verdiği bir araştırmada parametreyi,
  tahmin ediciyi ve gerçekleşmiş tahmini yazın.
  \item Yanlılığı $-1$ ve varyansı 5 olan bir tahmin edicinin MSE değerini bulun.
  \item A tahmin edicisinin yanlılığı 0, varyansı 12; B'nin yanlılığı 2,
  varyansı 4 ise MSE ölçütüne göre karşılaştırın.
  \item $3,5,7,9$ verisi için $n-1$ bölenli örneklem varyansını hesaplayın.
  \item Tutarlı bir tahmin edicinin küçük örneklemde neden yine de zayıf
  performans gösterebileceğini açıklayın.
  \item Örneklem ortalaması ile medyanı aykırı değerlere sağlamlık bakımından
  karşılaştırın.
  \item 30 kişilik grubun ortalaması 60, 70 kişilik grubun ortalaması 75 ise
  birleşik ortalamayı bulun.
  \item Standart hatası 0,08 olan bir oran tahminini dört ondalık basamakla
  raporlamanın neden sahte hassasiyet oluşturabileceğini açıklayın.
  \item Benzetim kullanarak bir tahmin edicinin yanlılık ve standart hatasını
  incelemek için izlenecek adımları yazın.
\end{enumerate}